\documentclass[10pt,a4paper]{amsart}
\usepackage{amsmath,amssymb,amsthm,graphicx}
\usepackage{amsmath,amssymb,amsthm,xfrac,xspace,hyperref,graphicx,comment}
\excludecomment{DefTralics}
\providecommand{\C}{}
\providecommand{\bibliofont}{\small}
\newcommand{\ie}{i.e.,\xspace}
\newcommand{\eg}{e.g.,\xspace}
\newcommand{\etc}{etc.\xspace}
\newcommand{\cf}{cf.\xspace}
\newcommand{\resp}{resp.\xspace}
\providecommand{\smaller}{\small}
\newcommand{\Subsection}{\subsection}
\newcommand{\Subsubsection}{\subsubsection}
\newcommand{\skpt}{\smallskip}
\emergencystretch=3em
\setcounter{section}{1}
\providecommand{\noopsort}[1]{}

\multlinegap0pt
\usepackage{mathrsfs}
\let\mathcal\mathscr
\newenvironment{enumeratea}
{\bgroup\def\theenumi{\alph{enumi}}\begin{enumerate}}
{\end{enumerate}\egroup}

\newcommand\mto{\mathchoice{\longmapsto}{\mapsto}{\mapsto}{\mapsto}}
\newcommand\hto{\mathrel{\lhook\joinrel\to}}
\newcommand\To[1]{\mathchoice{\xrightarrow{\textstyle\kern4pt#1\kern3pt}}{\stackrel{#1}{\longrightarrow}}{}{}}
\newcommand{\sbullet}{{\scriptscriptstyle\bullet}}
\newcommand{\tprod}{{\textstyle\prod}}
\newcommand{\rc}{\mathrm{c}}


\usepackage[all,cmtip]{xy}\let\labelstyle\textstyle

\theoremstyle{plain}
\newtheorem{theo}{Theorem}[section]
\newtheorem*{theo*}{Theorem}
\newtheorem{prop}[theo]{Proposition}
\newtheorem{lem}[theo]{Lemma}
\newtheorem{cor}[theo]{Corollary}

\theoremstyle{definition}
\newtheorem{defi}[theo]{Definition}
\newtheorem{rema}[theo]{Remark}
\newtheorem{conj}[theo]{Conjecture}

\DeclareFontFamily{U}{wncy}{}
\DeclareFontShape{U}{wncy}{m}{n}{%
<5>wncyr5%
<6>wncyr6%
<7>wncyr7%
<8>wncyr8%
<9>wncyr9%
<10>wncyr10%
<11>wncyr10%
<12>wncyr6%
<14>wncyr7%
<17>wncyr8%
<20>wncyr10%
<25>wncyr10}{}
\DeclareMathAlphabet{\cyrille}{U}{wncy}{m}{n}

\DeclareMathOperator \br {Br}
\DeclareMathOperator \cok {coker}
\DeclareMathOperator\Cone{Cone}
\DeclareMathOperator \Imm {Im}\let\Im\Imm
\DeclareMathOperator \Ker {Ker}
\DeclareMathOperator \pic {Pic}
\DeclareMathOperator \gal {Gal}
\DeclareMathOperator \spec {Spec}

\renewcommand\AA{\mathbf{A}}
\newcommand \ZZ {\mathbf{Z}}
\newcommand \QQ {\mathbf{Q}}
\newcommand \FF {\mathbf{F}}
\newcommand \LL {\mathbf{L}}
\newcommand \NN {\mathbf{N}}
\newcommand \RR {\mathbf{R}}
\newcommand \Gm {\mathbf{G}_m}\let\G\Gm

\newcommand\cF {\mathcal{F}}
\newcommand\cN {\mathcal{N}}
\newcommand\cO {\mathcal{O}}
\newcommand \cT {\mathcal{T}}
\newcommand \cC {\mathcal{C}}

\newcommand\Sha{\cyrille X}

\let\ov\overline

\newcommand \kbar {{\bar k}}
\newcommand \et {\textup{ét}}
\newcommand \nr {\mathrm{nr}}
\newcommand \tors {\mathrm{tors}}
\title{Artin-Mazur-Milne duality for bounded complexes of finite flat commutative group schemes over a curve}
\author{Cyril Demarche, David Harari}
\date{doi:10.5802/jep.129; CC BY 4.0}
\begin{document}
\maketitle
\noindent\textit{Editorial note.} The counted claim is the author Proposition AMM complex at source lines 1266-1272, and its proof is the proof environment that immediately follows it at source lines 1274-1311; no other proof in the paper is used for it.\par\medskip
\noindent\textit{Editorial note.} The section headings of the source at lines 196, 455 and 1233 are retained so that the printed numbering matches the published article, in which the counted statement is Proposition 4.1. The corpus record of the same paper's Theorem 4.9 (source label refin2, already counted as problem\_id 1933) confirms that the section entered here as number 4 is the published Section 4.\par\medskip
\noindent\textit{Editorial note.} The U-tori setup at source lines 455-470 is context only: the counted proposition concerns bounded complexes of finite flat commutative group schemes, and tori enter the paper only later.\par\medskip
\noindent\textit{Editorial note.} The inherited layout adapter sources/jep-2020-129-tori-duality/inherited-layout-adapter.tex redefines a command that is already defined under the amsart wrapper and fails to compile; the retained copy sources/jep-2020-129-tori-duality/tori-complex-layout-adapter.tex differs from it only in declaring that command with providecommand, and it is otherwise identical.\par\medskip
\noindent\textit{Editorial note.} The length-zero case of the induction is the cited Artin-Mazur-Milne theorem for one finite flat commutative group scheme; it is declared as a standard prerequisite, and every other step of the counted proof is retained in the body of this material.\par\medskip
\noindent\textit{Allowed standard prerequisites (editorial).}\begin{itemize}
\item Artin-Mazur-Milne duality for one finite flat commutative group scheme. For a finite flat commutative group scheme N over U there is a perfect pairing between the profinite group H\textasciicircum{}\{2-i\}\_c(U,N) and the discrete torsion group H\textasciicircum{}i(U, the Cartier dual of N), with values in H\textasciicircum{}3\_c(U,G\_m) = Q/Z.. Demarche-Harari, Artin-Mazur-Milne duality for fppf cohomology, Theorem 1.1, cited by the author at source line 1279 as the length-zero case.. Supplies only the length-zero case of the induction in the counted proof. The pairing for bounded complexes, its functoriality, the dualisation of the exact triangle and the five-lemma step are retained obligations of this task.
\item Value and finiteness of the top compact-support group. H\textasciicircum{}3\_c(U,G\_m) is isomorphic to Q/Z, and the topology on compact-support hypercohomology groups used here is the one already constructed in the retained source spans.. Milne, Arithmetic duality theorems, Proposition II.2.6, and Demarche-Harari, Proposition 2.1(4) and Proposition 3.5, cited at source lines 1275 and 444.. Identifies the target of the pairing and the profiniteness of the compact-support groups.
\item Homological algebra of bounded complexes. Bounded complexes of sheaves, their Hom-complexes with the sign conventions of the source, mapping cones, the cone-and-dual isomorphism, injective resolutions and the Godement resolution, together with the five lemma and induction on the length of a bounded complex.. The language of the statement and the mechanism of the counted proof.
\end{itemize}\medskip
\section{Compact support hypercohomology} \label{one}

Let $K$ be the function field of a smooth, projective, and geometrically
integral curve $X$ over a finite field $k$. Let $U$ be a non empty
Zariski open subset of $X$. Denote by $U^{(1)}$ the set of closed points
of $U$.

Let $\cC=(C_p)_{p \in \ZZ}$
be a bounded complex of fppf sheaves over $U$. In this text, we define the dual of $\cC$ to be the Hom-complex $\widehat{\cC}$ defined by
\[
\widehat{\cC}:= \underline{\textup{Hom}}^\sbullet(\cC, \Gm[1]),
\]
following the sign conventions in \cite[Tag 0A8H]{SP} or in \cite[X.5.1]{bkiac}.
Note that there is a functorial morphism of complexes
\[
\textup{Tot}(\cC \otimes \widehat{\cC}) \to \Gm[1]
\]
mapping an element $c \otimes \varphi\!\in\!C_p \otimes \textup{Hom}(C_q, A_n)$ to $0$ if $p\!\neq\!q$, and to $(-1)^{p (n-1)} \varphi(c)\!\in\!A_n$ if $p=q$.

With those conventions, if $\cC$ is concentrated in degree $0$, \ie $\cC = \cF$ with $\cF$ an fppf sheaf, then $\widehat{\cC}$ is the same as the Cartier dual $\cF^D:= \underline{\textup{Hom}}(\cF, \Gm)$ attached in degree $-1$, \ie $\widehat{\cC} = \cF^D[1]$ and the above pairing coincides with the obvious pairing $\cF \otimes \cF^D[1] \to \Gm[1]$ with no extra sign.

Note also that for any bounded complex $\cC$, we have a natural isomorphism of complexes $\widehat{\cC[1]} \To{\sim} \widehat{\cC}[-1]$, given by a sign $(-1)^{n+1}$ in degree $n$. And given a morphism $f: A \to B$ of bounded complexes, we have a natural isomorphism of complexes $\widehat{\textup{Cone}(f)} \To{\sim} \textup{Cone}(\widehat{f})[-1]$ such that the following diagram commutes\vspace*{-3pt}
\begin{equation} \label{cone and dual}
\begin{array}{c}
\xymatrix@R=6mm{
\widehat{A[1]} \ar[r] \ar[d]^\wr & \widehat{\textup{Cone}(f)} \ar[r] \ar[d]^\wr & \widehat{B} \ar[r]^{\widehat{f}} \ar@{=}[d] & \widehat{A} \ar@{=}[d] \\
\widehat{A}[-1] \ar[r] & \textup{Cone}(\widehat{f})[-1] \ar[r] & \widehat{B} \ar[r]^{\widehat{f}} & \widehat{A}.
}
\end{array}
\end{equation}

If $\cN$ is a commutative
group scheme (over $U$ or over $K$), its Cartier dual is
denoted $\cN^D$. The Pontryagin dual of a topological abelian group
$A$ (consisting of continuous homomorphism from $A$ to $\QQ/\ZZ$) is
denoted $A^*$. Unless explicitly specified, the topology used for
sheaves (\resp complex of sheaves) and cohomology (\resp hypercohomology) is the fppf topology.

For each closed point $v$ of $X$, the completion of $K$ at $v$
is denoted by $K_v$: it is a local field of characteristic $p$ with
finite residue field $\FF_v$ (observe the slight difference of notation with
\cite{DemH}, where $K_v$ stands for the henselization and $\widehat{K_v}$
for the completion). Denote by $\cO_v$ the ring of integers of $K_v$.
For every fppf sheaf $\cF$
over $U$ with generic fiber $F$, recall (\cite[Prop.\,2.1]{DemH})
the long exact sequence (where the piece of notation $v \not \in U$ means
that we consider all closed points of $X\setminus U$).
\begin{equation} \label{suppcomp}
\cdots\to H^i_\rc(U,\cF) \to H^i(U,\cF) \to \bigoplus_{v \not
\in U} H^i(K_v,F) \to H^{i+1}_\rc(U,\cF) \to\cdots
\end{equation}
There is also a long exact sequence
\begin{equation} \label{longex1}
\cdots\to H^i_\rc(U,\cF') \to H^i_\rc(U,\cF) \to
H^i_\rc(U,\cF'') \to H^{i+1}_\rc(U,\cF ') \to\cdots
\end{equation}
associated to every short exact sequence
\[
0 \to \cF' \to \cF \to \cF'' \to 0
\]
of fppf sheaves.

Let us now extend the construction of the groups $H^i_\rc(U,\dots)$ and \cite[Prop.\,2.1]{DemH} to the case of bounded complexes.
Let $\cC:= [\cdots\to \cF_i \to \cF_{i+1} \to\cdots]$ be a bounded complex of fppf sheaves over $U$.
Let $\cC \to I^\sbullet(\cC)$ be an injective resolution of the complex $\cC$, in the sense of \cite[Tag 013K]{SP}.
Following \cite[\S2]{DemH}, let $Z:= X \setminus U$ and
$Z':= \coprod_{v \in Z} \spec({K_v}) \To{i} U$.
Denote by $\cC_v$ and $I^\sbullet(\cC)_v$ their respective
pullbacks to $\spec K_v$, for $v \notin U$.

We now define $\Gamma_\rc(U, I^\sbullet(\cC))$ to be the following object in the category of complexes of abelian groups:
\[
\Gamma_\rc(U, I^\sbullet(\cC)):= \Cone\left(\Gamma(U, I^\sbullet(\cC)) \to \Gamma(Z', i^* I^\sbullet(\cC)) \right)[-1],
\]
and $H^r_\rc(U, \cC):= H^r(\Gamma_\rc(U, I^\sbullet(\cC)))$. We will also denote by $R\Gamma_\rc(U,\cC)$ the complex $\Gamma_\rc(U, I^\sbullet(\cC))$.
Similarly, one can define, for any closed point $v \in X$, complexes $\Gamma_v(\cO_v, \cC)$ computing fppf cohomology groups $H^r_v(\cO_v, \cC)$ over $\spec \cO_v$ with support in the closed point, as in \cite[before Lem.\,2.6]{DemH}.

As in \cite{DemH}, similar definitions could be made when $K$ is a number field (taking into account the real places), but in this article we will focus on the function field case. However, we will make remarks regularly throughout the text explaining similarities and differences appearing in the number field case.

We will need the analogue of \cite[Prop.\,2.1 \& 2.12]{DemH} for bounded complexes $\cC$: by construction, the first two points of
loc.\,cit., Prop.\,2.1 (\ie exact sequence \eqref{suppcomp} and~\eqref{longex1}) still hold for bounded complexes.

For a complex $\cC$ of finite flat group schemes,
let us now endow the groups $H^*_\rc(U, \cC)$ with a natural topology, compatible with the one defined in \cite{DemH} in the case where~$\cC$ is a finite flat group scheme.

Let $F$ be a local field (that is: a field complete for a discrete valuation
with finite residue field) and let $\cC:= [C_r \To{f_r} C_{r+1} \to \cdots \to C_s]$ be a bounded complex of finite commutative group schemes over $\spec F$, with $C_i$ in degree $i$. We assume that $F$ is of positive
characteristic $p$ (if $F$ is $p$-adic, then all
groups $H^r(F,\mathcal C)$ are finite by \cite[Cor.\,I.2.3]{MilADT}).

\begin{defi}
A morphism $f: G_1 \to G_2$
of topological groups is \emph{strict} if it is continuous, and
the restriction $f: G_1 \to f(G_1)$ is
an open map (where the topology on~$f(G_1)$ is induced by $G_2$). This is
equivalent to saying that $f$ induces an isomorphism of the topological
quotient $G_1/\ker f$ with the topological subspace $f(G_1) \subset G_2$.
\end{defi}

Let $A:= \ker(f_r)$, and let $\overline{\cC}:= \textup{Cone}(A[-r] \To{j} \cC)$. Then there is an exact triangle%
\begin{equation} \label{devissage complexe}
A[-r] \To{j} \cC \To{i} \overline{\cC} \To{p} A[1-r].
\end{equation}
In addition, we have a natural quasi-isomorphism $\varphi: \overline{\cC} \to \cC'$, where $\cC':= [C_{r+1}/\Im (f_r) \to C_{r+2} \to \cdots\to C_s]$ has a smaller length than $\cC$.

There is an alternative dévissage for the complex $\cC$, given by
the exact triangle:
\begin{equation} \label{devissage complexe bis}
\widetilde{\cC} \To{i'} \cC \To{p} C_r[-r] \To{\partial} \widetilde{\cC}[1],
\end{equation}
where $\widetilde{\cC}:= [C_{r+1} \To{f_{r+1}} C_{r+2} \to \cdots\to C_s]$.

Recall that for a finite and commutative $F$-group scheme $N$, the fppf groups
$H^i(F,N)$ are finite if $i \neq 1$ (\cite[\S III.6]{MilADT}) and they are
equipped with a locally compact topology for $i=1$ by \cite{Ces}.
By induction on the length of $\cC$, one deduces that if $C_i = 0$, then $H^{i+1}(F, \cC)$ is finite. In particular, with the previous notation, we get that $H^i(F, \cC)$ is finite if $i \leq r$ or $i \geq s+2$.

We now define a natural topology on $H^i(F, \cC)$ by induction on the length of $\cC$, such that any morphism of such complexes induces a strict map between hypercohomology groups. Using the dévissages given by \eqref{devissage complexe} and \eqref{devissage complexe bis}, one gets the following exact sequences\vspace*{-3pt}
\begin{multline} \label{devissage cohomo}
H^{i-1}(F, \cC') \to H^{i-r}(F,A) \To{f} H^i(F, \cC) \To{g} H^i(F, \cC')\\ \to H^{i+1-r}(F,A) \to \cdots\vspace*{-3pt}
\end{multline}
and\vspace*{-3pt}
\begin{multline} \label{devissage cohomo bis}
H^{i-r-1}(F, C_r) \to H^{i}(F,\widetilde{\cC}) \To{f'} H^i(F, \cC) \To{g'} H^{i-r}(F, C_r)\\ \to H^{i+1}(F,\widetilde{\cC}) \to \cdots,\vspace*{-3pt}
\end{multline}
where all the groups except $H^i(F, \cC)$ are endowed with a
natural topology via the induction hypothesis (observe that $C'_r=0$).
\begin{itemize}
\item Assume that $i = r+1$. Equip $\Im f\cong H^{i-r}(F,A)/
\Im H^{i-1}(F, \cC')$ with the quotient topology.
Then the two rightmost groups in exact sequence \eqref{devissage cohomo} are finite, and we can endow $H^{r+1}(F, \cC)$ with the topology such that
$\Im f$ is an open subgroup (see \cite[beginning of \S 3]{DemH}).
In other words it is the
finest topology such that $f$ is continuous. Then $f$ and $g$ are strict.
\item Assume that $i \neq r+1$. Then $H^{i-r}(F,C_r)$ is finite (and discrete), and using exact sequence \eqref{devissage cohomo bis} one can similarly
endow $H^i(F, \cC)$ with the finest topology making $f'$ continuous. Then both maps $f'$ and $g'$ are strict.
\end{itemize}
By construction, this topology is functorial in $\cC$, \ie if $\cC_1 \to \cC_2$ is a morphism of complexes, then the induced map $H^i(F, \cC_1) \to H^i(F, \cC_2)$ is strict. In addition, given a quasi-isomorphism $\cC_1 \to \cC_2$, the induced morphism on cohomology groups is a homeomorphism.

Let us now deal with the topology on the groups $H^*_\rc(U, \cC)$, where $\cC$ is a complex of finite flat commutative group schemes defined over $U$. Recall that we have an exact sequence analogous to \eqref{suppcomp}:
\begin{equation} \label{long exact seq compact supp}
H^{i-1}(U,\cC) \to \bigoplus_{v \notin U} H^{i-1}(K_v,\cC) \to H^{i}_\rc(U,\cC) \to H^{i}(U,\cC).
\end{equation}
We endow the groups $H^{i}(U,\!\cC)$ with the discrete topology, and the groups~$H^{i-1}(K_v,\!\cC)$ with the topology defined above.

\begin{lem} \label{closedimagecomplex}
For all $i$, the map $H^i(U, \cC) \to \bigoplus_{v \notin U} H^{i}(K_v,\cC)$ has discrete image.
\end{lem}

\begin{proof}
We prove the result by induction on the length of $\cC$, using the dévissages induced by the exact triangles \eqref{devissage complexe} and \eqref{devissage complexe bis}.
\begin{itemize}
\item Assume that $i = r+1$. Using the exact triangle \eqref{devissage complexe bis}, we get the following commutative diagram of long exact sequences of topological groups (where all maps are strict):\vspace*{-5pt}
\[
\xymatrix@R=6mm{
H^{r+1}(U, \widetilde{\cC}) \ar[r] \ar[d] & H^{r+1}(U, \cC) \ar[r] \ar[d] & H^1(U, C_r) \ar[d] \\
\bigoplus_{v \notin U} H^{r+1}(K_v, \widetilde{\cC}) \ar[r] & \bigoplus_{v \notin U} H^{r+1}(K_v, \cC) \ar[r] & \bigoplus_{v \notin U} H^1(K_v, C_r).
}
\]
The groups on the left hand side are finite, and by \cite[Lem.\,2.7]{cesnalms}, the image of the right hand side map is discrete in $\bigoplus_{v \notin U} H^1(K_v, C_r)$. Since $\bigoplus_{v \notin U} H^{r+1}(K_v, \cC)$ is Hausdorff, an easy topological argument implies that the image of the central vertical map is discrete.
\item Assume that $i \neq r+1$. Using the exact triangle \eqref{devissage complexe}, we get the following commutative diagram of long exact sequences of topological groups:\vspace*{-3pt}
\[
\xymatrix@R=6mm{
H^{i-r}(U, A) \ar[r] \ar[d] & H^{i}(U, \cC) \ar[r] \ar[d] & H^i(U, \cC') \ar[d] \\
\bigoplus_{v \notin U} H^{i-r}(K_v, A) \ar[r] & \bigoplus_{v \notin U} H^{i}(K_v, \cC) \ar[r] & \bigoplus_{v \notin U} H^i(K_v, \cC').
}
\]
The groups on the left hand side are finite, and by induction on the length of the complex, the image of the right hand side map is discrete in $\bigoplus_{v \notin U} H^1(K_v, \cC')$. A~similar topological argument as before implies that the central vertical map is discrete.\qedhere
\end{itemize}
\end{proof}

As a consequence of this Lemma, one can endow $H^i_\rc(U, \cC)$ with the following topology: we put the quotient topology on the group $\bigoplus_{v \notin U} H^{i}(K_v,\cC) / \Im H^i(U, \cC)$ (this topology is Hausdorff),
and since $H^i(U, \cC)$ is discrete, there is a unique topology on $H^i_\rc(U, \cC)$ such that the maps in the exact sequence \eqref{long exact seq compact supp} are strict.

\begin{lem} \label{lem topo profinite}
The topological group $H^i_\rc(U, \cC)$ is profinite.
\end{lem}

\begin{proof}
We prove this Lemma by induction on the length of the complex $\cC$.
By \cite[Prop.\,3.5]{DemH}, the lemma is proved when $\cC$ is a complex of length one, \ie concentrated in one given degree.

\pagebreak[2]
Given a complex $\cC$, consider the previous dévissages:
\begin{itemize}
\item Assume that $i = r+1$. Then the exact sequence \eqref{devissage cohomo} implies that the group $H^{r+1}_\rc(U, \cC)$ is an extension (the maps
being strict) of a (discrete) finite group by a profinite group (which is a quotient of $H^1_\rc(U, A)$ by a closed subgroup), hence $H^{r+1}_\rc(U, \cC)$ is profinite.
\item Assume that $i \neq r+1$. Then the exact sequence \eqref{devissage cohomo bis} implies that $H^i_\rc(U, \cC)$ is an extension of a finite (discrete) group by a profinite group (which is a quotient of $H^i_\rc(U, \widetilde{\cC})$ by a closed subgroup), hence it is profinite.\qedhere
\end{itemize}
\end{proof}

\section{Cohomology of tori and short complexes of tori} \label{onebis}

Let $U$ be a non empty
Zariski open subset of $X$. Recall that for every \emph{$U$-torus}~$\cT$ (in~the sense of \cite[IX, Déf.\,1.3]{sga3}),
there is a finite étale covering
(that can be taken to be connected and
Galois) $V$ of $U$ such that $\cT_V:=\cT \times_U V$
is \emph{split}, that is: isomorphic to some power $\G ^r$ ($r \in \NN$) of the
multiplicative group \hbox{(\cite[X, Th.\,5.16]{sga3})}. The \emph{group of characters}
$\widehat \cT$ of $\cT$ is a $U$-group scheme locally isomorphic to
$\ZZ^r$ for the étale topology, namely it is a torsion-free and finite
type $\gal(V/U)$-module.

Given a complex of $U$-tori $\cC=[\cT_1 \To{\rho} \cT_2]$
with generic fiber $C=[T_1 \To{\rho}T_2]$ over $K$,
where by convention $\cT_1$ is in degree $-1$ and $\cT_2$ in degree $0$, we

\section{Duality theorems in fppf cohomology} \label{two}

In order to state and prove duality results for the cohomology of complexes of fppf sheaves, we need to extend some constructions from \cite{DemH} to the context of bounded complexes.
Let $A$ and $B$ be two bounded complexes of fppf sheaves over $U$. Following \cite[XVII,\,4.2.9]{SGA4} or \cite[App.\,A]{friedsus}, one can consider the Godement resolutions $G(A)$ and $G(B)$ of $A$ and $B$. As in \cite[II.6.6]{God} or in \cite[App.\,A]{friedsus}, there is a natural commutative diagram of complexes
\[
\xymatrix{
A \otimes B \ar[rd] \ar[d] & \\
\textup{Tot}(G(A) \otimes G(B)) \ar[r] & G(A \otimes B).
}
\]
Following \cite[Proof of Lem.\,4.1]{DemH}, one gets a functorial morphism of complexes
\begin{equation} \label{cup complexes}
\textup{Tot}(\Gamma_\rc(U, G(A)) \otimes \Gamma(U,G(B))) \to \Gamma_\rc(U,G(A \otimes B))
\end{equation}
and a functorial pairing
\[
R\Gamma_\rc(U,A) \otimes^{\LL} R\Gamma(U,B) \to R\Gamma_\rc(U,A \otimes B).
\]

In particular, if $C$ is a bounded complex and $\widehat{C}:= \underline{\textup{Hom}}^\sbullet(C, \Gm[1])$ its dual, using the morphism $\textup{Tot}(C \otimes \widehat{C}) \to \Gm[1]$ from Section~\ref{one}, we get functorial pairings
\[\textup{Tot}(\Gamma_\rc(U, G(C)) \otimes \Gamma(U,G(\widehat{C}))) \to \Gamma_\rc(U,G(\Gm[1]))
\]
and
\begin{equation} \label{pairing derived}
R\Gamma_\rc(U,C) \otimes^{\LL} R\Gamma(U,\widehat{C}) \to R\Gamma_\rc(U,\Gm[1]).
\end{equation}

Following Section~\ref{one}, given a bounded complex $\cC$ of finite flat commutative group schemes over $U$, there is a natural topology on the abelian groups $H^i_\rc(U, \cC)$. This topology is profinite via Lemma \ref{lem topo profinite}, and considering $H^j(U, \widehat{\cC})$ as discrete torsion groups, the pairings
\[
H^i_\rc(U, \cC) \times H^j(U, \widehat{\cC}) \to H^{i+j+1}_\rc(U, \G)
\]
are continuous by the same argument as \cite[Lem.\,4.4]{DemH}.

\begin{prop} \label{AMM complex}
Let $\cC$ be a bounded complex of finite flat commutative group schemes over $U$.
For all $i \in \ZZ$, there is a perfect pairing between profinite and discrete torsion groups\vspace*{-3pt}
\[
H^{2-i}_\rc(U, \cC) \times H^i(U, \widehat{\cC}) \to H^3_\rc(U,\G) \cong \QQ/\ZZ.
\]
\end{prop}

\begin{proof}
The isomorphism $H^3_\rc(U,\G) \cong \QQ/\ZZ$ follows from \cite[Prop.\,II.2.6]{MilADT} and \cite[Prop.\,2.1\,(4)]{DemH}.

We now prove the proposition by induction on the length of the complex $\cC$.
\begin{itemize}
\item if $\cC$ is concentrated in a given degree $n$, then the proposition is a direct consequence of \cite[Th.\,1.1]{DemH}.
\item assume that $\cC:= [C_r \To{f_r} C_{r+11} \to \cdots\to C_s]$, with $C_i$ in degree $i$, has length $s-r \geq 1$.
Let $A:= \ker(f_r)$, and let $\overline{\cC}:= \textup{Cone}(A[-r] \To{j} \cC)$. Then there is an exact triangle\vspace*{-3pt}
\[
A[-r] \To{j} \cC \To{i} \overline{\cC} \To{p} A[1-r].
\]
We apply the functor $\widehat{\sbullet}$ to this exact triangle. Then, using \eqref{cone and dual}, we get that the natural triangle\vspace*{-3pt}
\[
\widehat{A[1-r]} \To{\widehat{p}} \widehat{\overline{\cC}} \To{\widehat{i}} \widehat{\cC} \To{\widehat{j}} \widehat{A[-r]}
\]
is exact.

Since the pairing between a complex and its dual is functorial, we get from the previous triangles a commutative diagram of topological groups with exact rows, where the vertical maps comes from the pairings \eqref{pairing derived}:\vspace*{-3pt}
\begin{multline} \label{devissage AMM1}
\begin{array}{c}
\xymatrix@=5mm{
H^{i-1}_\rc(U, \overline{\cC}) \ar[r] \ar[d] & H^{i-r}_\rc(U, A) \ar[r] \ar[d] & H^i_\rc(U, \cC)\ar[d] \\
H^{3-i}(U, \widehat{\overline{\cC}})^* \ar[r] & H^{3+r-i}(U, A^D)^* \ar[r] & H^{2-i}(U, \widehat{\cC})^*
}
\end{array}\\
\xymatrix@=5mm{
\mbox{}\ar[r] & H^{i}_\rc(U, \overline{\cC}) \ar[r] \ar[d] & H^{i+1-r}_\rc(U, A) \ar[d] \\
\mbox{}\ar[r] & H^{2-i}(U, \widehat{\overline{\cC}})^* \ar[r] & H^{2+r-i}(U, A^D)^*.
}
\end{multline}
The second line remains exact as the dual of an exact sequence of discrete groups.

Note that we have a quasi-isomorphism $\varphi: \overline{\cC} \to \cC'$, where $\cC':= [C_{r+1}/\Im(f_r) \to C_{r+2} \to \cdots\to C_s]$ has a smaller length than $\cC$, hence by induction, we can assume that the natural maps $H^r_\rc(U, \cC') \to H^{2-r}(U, \widehat{\cC'})^*$ are isomorphisms. Since all the $C_i$'s are finite flat group schemes, we see that the dual morphism $\widehat{\varphi}: \widehat{\cC'} \to \widehat{\overline{\cC}}$ is a quasi-isomorphism, hence by functoriality of the pairings, we deduce that the maps $H^r_\rc(U, \overline{\cC}) \to H^{2-r}(U, \widehat{\overline{\cC}})^*$ are isomorphisms too.

Hence in diagram \eqref{devissage AMM1}, all vertical morphisms, except perhaps the central one, are isomorphisms. Then the five lemma implies that the central morphism is an isomorphism.
\end{itemize}
By induction on the length of $\cC$, the proposition is proved.
\end{proof}

\begin{thebibliography}{99}
\bibitem{BD} M. Borovoi \& C. Demarche - “Manin obstruction to strong approximation for homogeneous spaces” , Comment. Math. Helv. 88 (2013) no. 1, p. 1-54
\bibitem{Bor} M. Borovoi - “The Brauer-Manin obstructions for homogeneous spaces with connected or abelian stabilizer” , J. reine angew. Math. 473 (1996), p. 181-194
\bibitem{bkiac} N. Bourbaki - Éléments de mathématique. Algèbre commutative. Chapitre 10 , Springer-Verlag, Berlin, 2007
\bibitem{Ces} K. Česnavičius - “Topology on cohomology of local fields” , Forum Math. Sigma 3 (2015), article ID e16, 55 pages
\bibitem{cesnalms} K. Česnavičius - “$p$-Selmer growth in extensions of degree $p$” , J. London Math. Soc. (2) 95 (2017) no. 3, p. 833-852
\bibitem{demlms} C. Demarche - “Le défaut d’approximation forte dans les groupes linéaires connexes” , Proc. London Math. Soc. (3) 102 (2011) no. 3, p. 563-597
\bibitem{demphd} C. Demarche - “Suites de Poitou-Tate pour les complexes de tores à deux termes” , Internat. Math. Res. Notices (2011) no. 1, p. 135-174, short version of “Théorèmes de dualité pour les complexes de tores”, available at https://webusers.imj-prg.fr/\textasciitilde{}cyril.demarche/
\bibitem{DemH} C. Demarche \& D. Harari - “Artin-Mazur-Milne duality for fppf cohomology” , Algebra Number Theory 13 (2019) no. 10, p. 2323-2357
\bibitem{friedsus} E. M. Friedlander \& A. Suslin - “The spectral sequence relating algebraic $K$-theory to motivic cohomology” , Ann. Sci. École Norm. Sup. (4) 35 (2002) no. 6, p. 773-875
\bibitem{gonz1} C. D. González-Avilés - “Arithmetic duality theorems for 1-motives over function fields” , J. reine angew. Math. 632 (2009), p. 203-231
\bibitem{gonz2} C. D. González-Avilés \& K.-S. Tan - “The generalized Cassels-Tate dual exact sequence for 1-motives” , Math. Res. Lett. 16 (2009) no. 5, p. 827-839
\bibitem{God} R. Godement - Topologie algébrique et théorie des faisceaux , Hermann, Paris, 1973
\bibitem{hss} D. Harari, C. Scheiderer \& T. Szamuely - “Weak approximation for tori over $p$-adic function fields” , Internat. Math. Res. Notices (2015) no. 10, p. 2751-2783
\bibitem{lang} S. Lang - “Algebraic groups over finite fields” , Amer. J. Math. 78 (1956), p. 555-563
\bibitem{MilEC} J. S. Milne - Étale cohomology , Princeton Math. Series , vol. 33 , Princeton University Press, Princeton, N.J., 1980
\bibitem{MilADT} J. S. Milne - Arithmetic duality theorems , BookSurge, LLC, Charleston, SC, 2006
\bibitem{nsw} J. Neukirch, A. Schmidt \& K. Wingberg - Cohomology of number fields , Grundlehren Math. Wiss. , vol. 323 , Springer-Verlag, Berlin, 2008
\bibitem{Ros} Z. Rosengarten - “Tamagawa numbers and other invariants of pseudo-reductive groups over global function fields” , 2018, to appear in Algebra Number Theory
\bibitem{San} J.-J. Sansuc - “Groupe de Brauer et arithmétique des groupes algébriques linéaires sur un corps de nombres” , J. reine angew. Math. 327 (1981), p. 12-80
\bibitem{col} J.-P. Serre - Corps locaux , Publications de l’Université de Nancago , vol. VIII , Hermann, Paris, 1968
\bibitem{SGA4} M. Artin, A. Grothendieck \& J.-L. Verdier - Théorie des topos et cohomologie étale des schémas. Tome 3 , Lect. Notes in Math. , vol. 269, 270 \& 305 , Springer-Verlag, Berlin-New York, 1973, Séminaire de Géométrie Algébrique du Bois-Marie 1963–1964 (SGA 4)
\bibitem{sga3} M. Demazure \& A. Grothendieck - Schémas en groupes (SGA 3). Tome I. Propriétés générales des schémas en groupes , Documents mathématiques , vol. 7 , Société Mathématique de France, Paris, 2011, article ID e16, Revised and annotated edition of the 1970 original
\bibitem{SP} Stacks Project Authors - “The Stacks Project” , 2019, https://stacks.math.columbia.edu
\end{thebibliography}
\end{document}
